% goldbach_G2.tex  ---  G2: 有限时间窗口（哥德巴赫 = [0,N) 时间平均，Polignac = 时间 -> 无穷 平均）
% 编译：pdflatex goldbach_G2.tex  （需要 CJKutf8 与 gbsn 字体；本机无 pdflatex，未编译，只做了静态配对检查）
% 记号与宏沿用 goldbach_props.tex（G1）。
\documentclass[11pt,a4paper]{article}
\usepackage{CJKutf8}
\usepackage{amsmath,amssymb,amsthm}
\usepackage{enumitem}
\usepackage{url}
\usepackage[margin=2.5cm]{geometry}

\newcommand{\Zhat}{\widehat{\mathbb Z}}
\newcommand{\Haar}{\mu}
\newcommand{\Sing}{\mathfrak S}
\newcommand{\Lsym}{\ensuremath{\mathsf{L}}}
\newcommand{\Rsym}{\ensuremath{\mathsf{R}}}
\newcommand{\status}[1]{\par\noindent\fbox{\textbf{状态：#1}}\par\smallskip}
% G2 新增记号
\newcommand{\Cnt}{\mathcal N}          % 窗口计数
\newcommand{\Tleg}{T}                   % 纯 Legendre 项数
\newcommand{\Tmix}{T^{*}}               % 混合界
\newcommand{\kmax}{k_{\max}}

\theoremstyle{plain}
\newtheorem{prop}{命题}[section]
\newtheorem{lem}[prop]{引理}
\newtheorem{cor}[prop]{推论}
\newtheorem{conj}[prop]{猜想}
\theoremstyle{definition}
\newtheorem{defn}[prop]{定义}
\newtheorem{rem}[prop]{注}
\renewcommand{\proofname}{证明}

\begin{document}
\begin{CJK}{UTF8}{gbsn}

\title{G2：有限时间窗口\\[2pt]
\large ——{}Legendre 误差的显式界、可控层 $\kmax(N)$ 与 Mertens 因子的来源}
\author{（G2 工作稿）}
\date{2026-11-02}
\maketitle

\begin{abstract}
沿用 G1（goldbach\_props.tex）的记号：$G_k^N=U_k\cap(N-U_k)\subset\Zhat$，里程计 $T(x)=x+1$。
哥德巴赫计数是 $0$ 的轨道在时间窗口 $[0,N)$ 内访问 $G_k^N$ 的次数，Polignac（固定层）读的是时间 $\to\infty$ 的平均
$\Haar(G_k^N)$；二者观测同一个 clopen 集。本稿给出窗口误差
$E_k(N)=\#\{0<n<N:n\in G_k^N\}-N\Haar(G_k^N)$ 的\emph{显式}界：
纯 Legendre 展开有 $\prod_{p\le p_k}(1+\nu_p)=3^{k-\omega_k(N)}2^{\omega_k(N)}$ 项（不是 $2^k$ 项），
与周期性结合后 $|E_k(N)|\le\prod_{p\le p_k}\min(p-\nu_p,1+\nu_p)\le2\cdot3^{k-2}$。
由此得到：Legendre 可控的最大层数为 $\kmax(N)=\bigl(\log(\varepsilon N)-2\log\log\log N+O(1)\bigr)/\log3$，
对应 $\log P_k/\log N=(1+o(1))\log\log N/\log3$——即使 $P_k\gg N$（窗口短于周期）也可控，
真正的约束是剩余类的\emph{个数} $3^k$ 而不是周期 $P_k$。
在 $k=\pi(\sqrt N)$ 处 $P_k=e^{(1+o(1))\sqrt N}$，Legendre 界超过主项 $e^{c\sqrt N/\log N}$ 倍；
但\emph{真实}误差只是主项的一个常数比例，在 Hardy--Littlewood 猜想下趋于 $1-e^{2\gamma}/4\approx20.7\%$——
这正是 Mertens 因子 $4e^{-2\gamma}=(e^\gamma\omega(2))^{-2}$（$\omega$ 为 Buchstab 函数）。
最后与 Codex 的 D2/D3 数值对照。

\textbf{定位声明。} 本稿全部严格结论都是初等筛法（Legendre、Mertens、素数定理）的直接推论或对经典结果
（Buchstab、筛法基本引理、Selberg 上界筛）的引用；\emph{不}对哥德巴赫或 Polignac 猜想提供任何新证据。
数值只作说明，不用于证明任何渐近陈述。
\end{abstract}

\section*{状态总表}
\begin{center}
\begin{tabular}{p{2.6cm}p{7.6cm}p{4.2cm}}
\hline
条目 & 内容 & 状态\\
\hline
引理 \ref{lem:local} & $G_k^N$ 的局部结构：$p\nmid N$ 删 2 类，$p\mid N$ 删 1 类；Legendre 项数 & 已证\\
命题 \ref{prop:legendre} & 窗口误差的显式界 $|E|\le\Tmix_k(N)\le2\cdot3^{k-2}$ & 已证\\
推论 \ref{cor:relative} & 相对误差的显式界 $\beta_k(N)$，常数 $c_*=e^{2\gamma}/(18C_2)$ & 已证（显式常数依赖引用的 Rosser--Schoenfeld 不等式）\\
命题 \ref{prop:kmax} & $\kmax(N,\varepsilon)$ 的定义、阶、$\theta$ 的允许范围与其锐性 & 已证\\
命题 \ref{prop:top} & $k=\pi(\sqrt N)$：$P_k=e^{(1+o(1))\sqrt N}$，Legendre 界/主项 $=e^{\Theta(\sqrt N/\log N)}$ & 已证\\
命题 \ref{prop:toperr} & 顶层真实误差：(a) 与 $r(N)$ 的精确关系；(b) $O(1)$ 倍主项；(c) 极限 $e^{2\gamma}/4-1$ & (a) 已证；(b) 部分（引用）；(c) 猜想（$\Leftrightarrow$ HL）\\
命题 \ref{prop:buchstab} & 一维 Buchstab：$\Phi(x,x^{1/u})/(x\prod(1-\frac1p))\to e^\gamma\omega(u)$ & 已证（经典，引用）\\
命题 \ref{prop:FL} & 筛法基本引理给出的可控层 $p_k\le N^{c/\log(1/\varepsilon)}$ & 部分（引用，常数不显式）\\
猜想 \ref{conj:2dim} & 二维 Buchstab：$\Cnt_k(N)/(N\Haar)\to(e^\gamma\omega(s))^2$ & 猜想\\
注 \ref{rem:mertenssource} & Mertens 因子的来源及与 G1 注 rem:mertensrate 的衔接 & 解释（所用恒等式已证）\\
\S\ref{sec:D2} & 与 D2/D3 的对照 & 数值说明，非证明\\
\hline
\end{tabular}
\end{center}

%=====================================================================
\section{记号与 $G_k^N$ 的局部结构}\label{sec:notation}

记号同 G1：$p_1=2<p_2=3<\cdots$，$P_k=\prod_{i\le k}p_i$，$U_k=\{x\in\Zhat:x_p\ne0,\ p\le p_k\}$，
$G_k^N=U_k\cap(N-U_k)$，$\Haar$ 为 $\Zhat$ 上的 Haar 概率测度，$\gamma$ 为 Euler 常数，
$C_2=\prod_{p\ge3}(1-(p-1)^{-2})=0.6601618158\ldots$，$\Sing(N)$ 为 Hardy--Littlewood 奇异级数，
$R_k(N)=\Haar(G_k^N)/\Haar(U_k)^2$（G1 命题 prop:ratio）。此外：
\begin{itemize}[leftmargin=*]
\item $\nu_p(N):=1$（$p\mid N$），$\nu_p(N):=2$（$p\nmid N$）；$\omega_k(N):=\#\{p\le p_k:p\mid N\}$；
对无平方因子 $d\mid P_k$，$\nu(d):=\prod_{p\mid d}\nu_p(N)$。
\item 窗口计数与误差（与 D2 的定义逐字一致）：
\[
\Cnt_k(N):=\#\{0<n<N:\ n\in G_k^N\},\qquad
E_k(N):=\Cnt_k(N)-N\Haar(G_k^N),\qquad
\rho_k(N):=\frac{|E_k(N)|}{N\Haar(G_k^N)} .
\]
更一般地，对整数 $x_0$ 与 $L\ge0$，$E_k(N;x_0,L):=\#\{x_0\le n<x_0+L:n\in G_k^N\}-L\Haar(G_k^N)$。
\item 由于 $0\notin U_k$，有 $0\notin G_k^N$，$N\notin G_k^N$（因 $N-N=0$）。故
$\Cnt_k(N)=\#\{0\le n<N:n\in G_k^N\}$，即 $E_k(N)=E_k(N;0,N)$：\emph{窗口恰为 $N$ 个连续整数}，
这正是里程计 $T$ 从 $0$ 出发、时间 $0,1,\dots,N-1$ 的轨道段（G1 命题 prop:odometer）。
\end{itemize}

\begin{lem}[局部结构与 Legendre 项数]\label{lem:local}
设 $N\in\mathbb Z$，$k\ge1$。
\begin{enumerate}[label=(\alph*)]
\item $n\in G_k^N\iff$ 对每个 $p\le p_k$，$n\bmod p\notin A_p:=\{0,\ N\bmod p\}$；$\#A_p=\nu_p(N)$，
即 $p\nmid N$ 时删去 2 个剩余类，$p\mid N$ 时删去 1 个。
\item $\Haar(G_k^N)=\prod_{p\le p_k}\bigl(1-\nu_p(N)/p\bigr)$，$\#(G_k^N\bmod P_k)=\prod_{p\le p_k}(p-\nu_p(N))$。
$N$ 奇时（$p=2$ 处 $\nu_2=2$）$G_k^N=\varnothing$；$N$ 偶时 $\nu_2=1$。
\item（Legendre 展开）
\[
\mathbf 1_{G_k^N}(n)=\sum_{d\mid P_k}\mu(d)\,\mathbf 1\bigl[n\bmod p\in A_p\ \forall p\mid d\bigr],
\]
其中第 $d$ 项的条件由 CRT 是模 $d$ 的 $\nu(d)$ 个剩余类之并。按“（$d$，剩余类）”计，展开共有
\[
\Tleg_k(N):=\sum_{d\mid P_k}\nu(d)=\prod_{p\le p_k}\bigl(1+\nu_p(N)\bigr)=3^{\,k-\omega_k(N)}\,2^{\,\omega_k(N)}
\]
项（含 $d=1$）。特别地，对偶数 $N$，$2^{k}\cdot(3/2)^{\,k-\omega_k(N)}=\Tleg_k(N)\le2\cdot3^{k-1}$。
\end{enumerate}
\end{lem}
\status{已证}
\begin{proof}
(a) $n\in U_k\iff n_p\neq0$；$n\in N-U_k\iff (N-n)_p\ne0\iff n_p\ne N_p$。$N_p=0$ 时两条件重合。
(b) CRT 下各 $x_p$ 独立均匀。(c) $\mathbf 1_{G}=\prod_{p\le p_k}(1-\mathbf 1_{A_p}(n\bmod p))$ 展开；
$\prod_{p\mid d}\mathbf 1_{A_p}$ 是 $\prod_{p\mid d}\#A_p=\nu(d)$ 个模 $d$ 类之并；
$\sum_{d\mid P_k}\nu(d)=\prod_p(1+\nu_p)$ 因 $\nu$ 积性。偶数 $N$ 有 $p=2\mid N$，故 $\omega_k(N)\ge1$。
\end{proof}

\begin{rem}[关于任务中“$O(2^{\omega(P_k)})$ 级误差”]\label{rem:twoomega}
$2^{\omega(P_k)}=2^k$ 是无平方因子因子 $d\mid P_k$ 的\emph{个数}；但每个 $d$ 携带 $\nu(d)$ 个剩余类，
每个剩余类在窗口中产生一个 $(-1,1)$ 内的舍入误差，所以 Legendre 论证给出的是 $3^{k-\omega_k(N)}2^{\omega_k(N)}$ 量级，
而不是 $2^k$。一个模 $d$ 的 $\nu(d)$ 类之并在长度 $L$ 的区间中的计数误差一般只能界为 $\min(\nu(d),d-\nu(d))$，
不能界为 $1$，故 $|E|\le C\cdot2^k$ 型界\emph{不能}由 Legendre 论证得到。我们既未证明也未否证 $|E_k(N;x_0,L)|\ll2^k$；
作为说明（g2\_scratch/ 中暴力枚举，$k\le6$、全部偶 $N\bmod P_k$、全部窗口 $[a,b)$，$0\le a<b\le P_k$）：最大偏差依次为
$0.50,\,1.00,\,1.80,\,3.86,\,7.52,\,16.52$，均 $<2^k$，但这对 $k\to\infty$ 不说明任何问题。
下面的显式界一律用 $3^{\cdot}2^{\cdot}$ 型。
\end{rem}

%=====================================================================
\section{命题 2.1：显式 Legendre 误差与可控层}\label{sec:legendre}

\begin{prop}[窗口误差的显式界]\label{prop:legendre}
设 $N\in\mathbb Z$，$k\ge1$，$x_0\in\mathbb Z$，$L\ge0$。则
\begin{enumerate}[label=(\roman*)]
\item（纯 Legendre）$|E_k(N;x_0,L)|\le\Tleg_k(N)-1$；
\item（周期性）$|E_k(N;x_0,L)|\le\#(G_k^N\bmod P_k)=P_k\Haar(G_k^N)$；
\item（混合）$|E_k(N;x_0,L)|\le\Tmix_k(N):=\prod_{p\le p_k}\min\bigl(p-\nu_p(N),\,1+\nu_p(N)\bigr)$。
\end{enumerate}
对偶数 $N$ 与 $k\ge2$：$\Tmix_k(N)=m_3(N)\prod_{5\le p\le p_k}(1+\nu_p(N))$，其中 $m_3(N)=1$（$3\nmid N$）或 $2$（$3\mid N$）；
因此
\[
|E_k(N)|\le\Tmix_k(N)=m_3(N)\,3^{\,k-2-\omega'_k(N)}\,2^{\,\omega'_k(N)}\le2\cdot3^{k-2},\qquad
\omega'_k(N):=\#\{5\le p\le p_k:\ p\mid N\}.
\]
$N$ 奇时 $E_k\equiv0$（$\Tmix_k=0$）。
\end{prop}
\status{已证}
\begin{proof}
基本事实：模 $q$ 的一个剩余类在任意 $L$ 个连续整数中出现 $L/q+\delta$ 次，$\delta\in(-1,1)$。

(iii) 令 $S:=\{p\le p_k:\ p-\nu_p\le1+\nu_p\}$，$Q:=\prod_{p\in S}p$，$B\subset\mathbb Z/Q$ 为在 $S$ 中各素数处都允许的类，
$\#B=\prod_{p\in S}(p-\nu_p)$。记 $P'=P_k/Q$。则
\[
\mathbf 1_{G_k^N}(n)=\mathbf 1_B(n\bmod Q)\prod_{p\mid P'}\bigl(1-\mathbf 1_{A_p}(n\bmod p)\bigr)
=\sum_{d\mid P'}\mu(d)\,\mathbf 1_B(n\bmod Q)\,\mathbf 1[n\bmod p\in A_p\ \forall p\mid d].
\]
由 CRT（$\gcd(Q,d)=1$），第 $d$ 项是模 $Qd$ 的 $\#B\cdot\nu(d)$ 个类之并，在窗口中的计数为
$L\,\#B\,\nu(d)/(Qd)+\delta_d$，$|\delta_d|\le\#B\,\nu(d)$。主项之和为
$L\frac{\#B}{Q}\prod_{p\mid P'}(1-\frac{\nu_p}{p})=L\Haar(G_k^N)$（引理 \ref{lem:local}(b)）。故
$|E|\le\#B\sum_{d\mid P'}\nu(d)=\prod_{p\in S}(p-\nu_p)\prod_{p\notin S}(1+\nu_p)=\prod_p\min(p-\nu_p,1+\nu_p)$。

(i) 取 $S=\varnothing$：同上，但 $d=1$ 项的计数恰为 $L$（无误差），故 $|E|\le\Tleg_k-1$。

(ii) 写 $L=qP_k+r$，$0\le r<P_k$。完整周期贡献 $q\#G$ 与主项 $qP_k\Haar(G)=q\#G$ 精确相消；
余下长度 $r$ 的段中计数 $\in[0,\#G]$，$r\Haar(G)\in[0,\#G)$，差的绝对值 $\le\#G$。

偶数 $N$：$p=2$ 时 $\min(1,2)=1$；$p=3$ 时 $\min(1,3)=1$（$3\nmid N$）或 $\min(2,2)=2$（$3\mid N$）；
$p\ge5$ 时 $p-\nu_p\ge3\ge1+\nu_p$，取 $1+\nu_p\in\{2,3\}$。$N$ 奇时 $p=2$ 处 $p-\nu_2=0$。
\end{proof}

\begin{rem}[数值自洽（说明性）]
g2\_scratch/check\_g2.py 对 $N<200$ 及若干特殊 $N$、$k\le5$、多种 $(x_0,L)$ 暴力核对了 (i)(iii)，
所见最大 $|E|/\Tmix_k=0.83$；并对 results/D2.json 的全部 9871 个 $(N,k)$ 层（字段 \texttt{layers[].E}）核对
$|E_k(N)|\le\Tmix_k(N)$，无一违反。这只是自洽检查，命题 \ref{prop:legendre} 的证明不依赖它。
\end{rem}

\begin{cor}[相对误差的显式界]\label{cor:relative}
对偶数 $N$、$k\ge2$：
\[
\rho_k(N)\le\beta_k(N):=\frac{\Tmix_k(N)}{N\Haar(G_k^N)}=\frac6N\prod_{5\le p\le p_k}f_p(N),\qquad
f_p(N)=\begin{cases}\dfrac{2p}{p-1},&p\mid N,\\[6pt]\dfrac{3p}{p-2},&p\nmid N,\end{cases}
\]
从而 $\beta_k(N)\le\beta^\sharp_k(N):=\dfrac6N\,3^{k-2}\prod_{5\le p\le p_k}\Bigl(1-\dfrac2p\Bigr)^{-1}$（与 $N$ 的算术无关），并且
\[
\beta^\sharp_k(N)\le c_*\,\frac{3^k\,(\log p_k)^2\bigl(1+(\log p_k)^{-2}\bigr)^2}{N},\qquad
c_*:=\frac{e^{2\gamma}}{18\,C_2}=0.26695\ldots
\]
\end{cor}
\status{已证（最后一个不等式依赖 Rosser--Schoenfeld 显式 Mertens 界 \cite{RS62}，见证明）}
\begin{proof}
$\beta_k=\frac1N\prod_{p\le p_k}\frac{p\min(p-\nu_p,1+\nu_p)}{p-\nu_p}$；$p=2$ 因子为 $2$，$p=3$ 因子在两种情形都为 $3$，
$p\ge5$ 因子为 $p(1+\nu_p)/(p-\nu_p)=f_p$。$3p/(p-2)>2p/(p-1)$ 给出 $\beta_k\le\beta_k^\sharp$。
对最后一式：$(1-\frac2p)^{-1}=(1-\frac1p)^{-2}(1-\frac1{(p-1)^2})^{-1}$；
$\prod_{p\ge5}(1-\frac1{(p-1)^2})^{-1}=\frac{3/4}{C_2}$；$\prod_{5\le p\le y}(1-\frac1p)^{-2}=\frac19\bigl(\prod_{p\le y}\frac p{p-1}\bigr)^2$；
再用 \cite[Thm.~8, (3.30)]{RS62}：$\prod_{p\le y}\frac p{p-1}<e^\gamma\log y\,(1+(\log y)^{-2})$（$y>1$；条件以原文为准，
本稿仅对 $y\le2\times10^5$ 的素数点数值抽查过，最大比值 $0.9997$）。合并：$\frac6N3^{k-2}\cdot\frac{3/4}{C_2}\cdot\frac19\cdot e^{2\gamma}(\log p_k)^2(1+\cdots)^2
=\frac{e^{2\gamma}}{18C_2}\frac{3^k(\log p_k)^2(1+\cdots)^2}{N}$。
\end{proof}

\begin{prop}[Legendre 可控的最大层 $\kmax(N)$ 与 $\theta$ 的允许范围]\label{prop:kmax}
对 $\varepsilon\in(0,1)$ 与偶数 $N$，定义
\[
\kmax(N,\varepsilon):=\max\{k\ge2:\ \beta^\sharp_k(N)\le\varepsilon\}
\quad(\text{集合为空时记 }1),
\]
以及依赖 $N$ 的版本 $\kmax^{(N)}(N,\varepsilon):=\max\{k\ge2:\beta_k(N)\le\varepsilon\}\ge\kmax(N,\varepsilon)$。
\begin{enumerate}[label=(\alph*)]
\item（可控性）$\beta_k^\sharp$、$\beta_k(N)$ 关于 $k$ 严格递增；对一切 $2\le k\le\kmax(N,\varepsilon)$，$\rho_k(N)\le\varepsilon$。
充分条件：$c_*3^k(\log p_k)^2(1+(\log p_k)^{-2})^2\le\varepsilon N$。
\item（阶）对固定 $\varepsilon$（更一般地，对 $N^{-1/2}\le\varepsilon<1$ 一致），
\[
\kmax(N,\varepsilon)=\Bigl\lfloor\frac{\log(\varepsilon N)-2\log\log\log N+\log(1/c_*)}{\log3}+o(1)\Bigr\rfloor
=\frac{\log(\varepsilon N)-2\log\log\log N}{\log 3}+O(1),
\]
$\log(1/c_*)=1.3207\ldots$。相应地 $p_{\kmax}\sim\frac{\log N\,\log\log N}{\log3}$，
\[
x_{\max}(N):=\frac{\log P_{\kmax}}{\log N}=(1+o(1))\frac{\log\log N}{\log 3}\longrightarrow\infty .
\]
\item（$\theta$ 的允许范围）若 $\theta=\theta(N)\le(1-\delta)\frac{\log\log N}{\log3}$（$\delta>0$ 固定），则凡 $P_k\le N^\theta$ 者
$\rho_k(N)\le\beta_k^\sharp(N)\le N^{-\delta+o(1)}$；特别地，对\emph{任意固定} $\theta>0$，$P_k\le N^{\theta}$ 蕴含
$\rho_k(N)\le N^{-1+O_\theta(1/\log\log N)}$。
\item（锐性）反之，对\emph{每个}偶数 $N$，若 $k\ge(1+\delta)\frac{\log N}{\log3}$（等价地 $\log P_k/\log N\ge(1+\delta+o(1))\frac{\log\log N}{\log3}$），
则 $\beta_k(N)\ge N^{\delta-o(1)}\to\infty$：混合界本身失效。
\end{enumerate}
\end{prop}
\status{已证（(b)--(d) 用 Mertens 第三定理与素数定理 $\vartheta(y)\sim y$，仅用于渐近，不需显式常数）}
\begin{proof}
(a) 由推论 \ref{cor:relative} 及 $f_p>1$。
(b) 由 Mertens 第三定理，$\prod_{3\le p\le y}(1-\frac2p)=4C_2(y)\prod_{p\le y}(1-\frac1p)^2\sim4C_2e^{-2\gamma}/\log^2y$，
其中 $C_2(y)=\prod_{3\le p\le y}(1-(p-1)^{-2})$；又 $1-\frac23=\frac13$，故
$\prod_{5\le p\le y}(1-\frac2p)^{-1}\sim\frac{e^{2\gamma}}{12C_2}\log^2y$。于是
\[
\log\beta_k^\sharp=-\log N+k\log3+2\log\log p_k+\log c_*+o(1).
\]
相邻两层之差为 $\log\frac{3p_{k+1}}{p_{k+1}-2}=\log3+o(1)$，故 $\kmax$ 为方程 $\log\beta^\sharp=\log\varepsilon$ 实根的整数部分（误差 $o(1)$）。
在 $k\asymp\log N$ 处 $p_k\sim k\log k$，$\log\log p_k=\log\log\log N+o(1)$，得第一式。
$P_k=e^{\vartheta(p_k)}$，$\vartheta(p_k)\sim p_k\sim\kmax\log\kmax\sim\frac{\log N}{\log3}\log\log N$。
(c) 由单调性，只需考虑最大的 $k$ 使 $P_k\le N^\theta$，且可设 $\theta\ge1$。$\vartheta(p_k)\le\theta\log N$ 给出
$p_k\le(1+o(1))\theta\log N$，$k=\pi(p_k)\le(1+o(1))\frac{\theta\log N}{\log(\theta\log N)}\le(1+o(1))(1-\delta)\frac{\log N}{\log3}$。
于是 $3^k\le N^{(1-\delta)(1+o(1))}$，而 $(\log p_k)^2=N^{o(1)}$，代入推论 \ref{cor:relative}。固定 $\theta$ 时
$k\le(1+o(1))\theta\log N/\log\log N$，$3^k=N^{O(\theta/\log\log N)}$。
(d) $f_p\ge3$（$p\nmid N$），$f_p\ge2$（$p\mid N$），故 $\beta_k(N)\ge\frac6N3^{k-2}(2/3)^{\omega(N)}$；
而 $\omega(N)\le(1+o(1))\log N/\log\log N$，$(2/3)^{\omega(N)}=N^{-o(1)}$。
\end{proof}

\begin{rem}[数值：由定义直接计算的 $\kmax(10^m,0.1)$]\label{rem:kmaxtable}
下表是对上述\emph{定义}的确定性计算（不是实验数据）：
\begin{center}
\begin{tabular}{c|ccccccccc}
\hline
$N$ & $10^4$ & $10^5$ & $10^6$ & $10^7$ & $10^8$ & $10^9$ & $10^{10}$ & $10^{11}$ & $10^{12}$\\
\hline
$\kmax(N,0.1)$ & 5 & 7 & 9 & 11 & 13 & 15 & 17 & 19 & 21\\
$p_{\kmax}$ & 11 & 17 & 23 & 31 & 41 & 47 & 59 & 67 & 73\\
$x_{\max}=\log P_{\kmax}/\log N$ & 0.84 & 1.14 & 1.39 & 1.61 & 1.81 & 1.98 & 2.13 & 2.26 & 2.38\\
$\log\log N/\log 3$ & 2.02 & 2.22 & 2.39 & 2.53 & 2.65 & 2.76 & 2.86 & 2.94 & 3.02\\
\hline
\end{tabular}
\end{center}
渐近式 $x_{\max}\sim\log\log N/\log3$ 收敛很慢（$\log\log p_k$ 与 $\log\log\log N$ 之差在此范围内仍约 $0.2$--$0.3$）；
表中数字只说明量级，不用来“验证”渐近式。要点：$N\ge10^5$ 起 $x_{\max}>1$，即 Legendre 在 $P_k>N$（窗口短于周期）时仍可控。
\end{rem}

\begin{rem}[“窗口短于周期”不是 Legendre 的瓶颈]
任务的提法“$P_k\le N^\theta$”暗示需要 $\theta<1$（窗口至少一个周期）。命题 \ref{prop:kmax}(c)(d) 表明真正的阈值是
剩余类展开的项数 $\Tmix_k\approx3^k$ 对比 $N$，即 $k\approx\log N/\log3$；用周期表示，阈值为
$\theta^*(N)=(1+o(1))\log\log N/\log3$，它随 $N$ 缓慢趋于无穷。$\theta<1$ 只是一个（过于保守的）充分条件。
\end{rem}

%=====================================================================
\section{顶层 $k=\pi(\sqrt N)$：窗口太短}\label{sec:top}

\begin{prop}[顶层的周期与 Legendre 界]\label{prop:top}
设 $N$ 偶，$k=\pi(\sqrt N)$。则
\begin{enumerate}[label=(\alph*)]
\item $\log P_k=\vartheta(\sqrt N)=(1+o(1))\sqrt N$，即 $P_k=e^{(1+o(1))\sqrt N}\gg N$；窗口只覆盖一个周期的
$N/P_k=e^{-(1+o(1))\sqrt N}$ 部分。
\item $(2\log2+o(1))\dfrac{\sqrt N}{\log N}\le\log\beta_k(N)\le(2\log3+o(1))\dfrac{\sqrt N}{\log N}$：
混合 Legendre 界超过主项 $e^{\Theta(\sqrt N/\log N)}$ 倍；纯 Legendre 界 $\Tleg_k-1$ 与周期界 $P_k\Haar$ 更差。
\end{enumerate}
\end{prop}
\status{已证（素数定理）}
\begin{proof}
(a) 素数定理 $\vartheta(y)\sim y$。(b) 上界：$\beta_k\le\beta_k^\sharp$，且 $\log\prod_{5\le p\le\sqrt N}(1-\frac2p)^{-1}=O(\log\log N)$，
$k\log3=(2\log3+o(1))\sqrt N/\log N$。下界：$f_p>2$，故 $\beta_k\ge\frac6N2^{k-2}$，$k\log2=(2\log2+o(1))\sqrt N/\log N\gg\log N$。
\end{proof}

\begin{rem}[动力学解读：时间 $N$ 只“看见”模 $\le N$ 的结构]\label{rem:level}
里程计轨道 $0,1,\dots,N-1$ 对每个模 $d$ 的每个剩余类恰好命中 $N/d+\delta$（$|\delta|<1$）次，这是\emph{精确}的；
但当 $d>N$ 时每类只命中 $0$ 或 $1$ 次，而 Haar 主项 $N/d<1$ 是“虚构”的。$G_k^N$ 是模 $P_k$ 的
$\prod(p-\nu_p)\approx e^{\sqrt N}/\log^2N$ 个类之并，Legendre 展开中绝大多数 $d\mid P_k$ 满足 $d>N$，其舍入误差不再相消。
用筛法语言：时间窗口 $[0,N)$ 提供的分布水平（level）为 $D\approx N$，筛到 $z=p_k$；
参数 $s:=\log D/\log z=\log N/\log p_k$。哥德巴赫的顶层是 $s=2$。
Polignac（G1 命题 prop:window 之后的注）在\emph{固定} $k$ 下读时间 $\to\infty$ 平均，得到精确的 $\Haar(G_k^N)$
（命题 \ref{prop:legendre}(ii)：$|E_k(N;0,L)|/L\le P_k\Haar/L\to0$）；而真正的孪生/间隔 $N$ 素数计数也要求 $k$ 随窗口 $X$ 增长到
$p_k\approx\sqrt X$，同样处于 $s=2$。二者共享同一个困难。
\end{rem}

\begin{prop}[顶层的真实误差]\label{prop:toperr}
设 $N\ge4$ 偶，$k=\pi(\sqrt N)$，$r(N)$ 为有序哥德巴赫表示数。
\begin{enumerate}[label=(\alph*)]
\item $-2\pi(\sqrt N)\le\Cnt_k(N)-r(N)\le2$。因此
\[
\frac{E_k(N)}{N\Haar(G_k^N)}=\frac{r(N)}{N\Haar(G_k^N)}-1+O\!\Bigl(\frac{\log N}{\sqrt N}\Bigr).
\]
\item $-1\le E_k(N)/(N\Haar(G_k^N))\le2e^{2\gamma}-1+o(1)\approx5.34$。
\item 下面两件事等价：(i) Hardy--Littlewood 猜想 $r(N)\sim\Sing(N)N/(\log N)^2$（偶数 $N\to\infty$）；
(ii) $E_k(N)/(N\Haar(G_k^N))\to e^{2\gamma}/4-1=-0.20694\ldots$。
\end{enumerate}
\end{prop}
\status{(a) 已证；(b) 部分（上界引用 Selberg 上界筛 \cite{HR74}）；(c) 中的等价已证，而 (i)/(ii) 本身是猜想}
\begin{proof}
(a) 若 $1<n\le\sqrt N$ 且 $n\in G_k^N$，则 $n$ 的最小素因子 $\le n\le\sqrt N$，即 $\le p_k$，矛盾；故 $[0,\sqrt N]$ 中至多 $n=1$。
同理（对 $N-n$）$[N-\sqrt N,N)$ 中至多 $n=N-1$。中段 $\sqrt N<n<N-\sqrt N$ 的计数由 G1 命题 prop:window 介于
$r(N)-2\pi(\sqrt N)$ 与 $r(N)$ 之间。误差项：$N\Haar(G_k^N)=N\Haar(U_k)^2R_k(N)\ge N\Haar(U_k)^2\cdot2C_2\gg N/\log^2N$
（Mertens；$R_k(N)\ge2C_2(k)\ge2C_2$，G1 命题 prop:ratio(b)），而 $\pi(\sqrt N)\ll\sqrt N/\log N$。

(b) 下界因 $\Cnt\ge0$。上界：$\Cnt_k(N)$ 是满足 $n,N-n$ 都没有 $\le\sqrt N$ 的素因子的 $n$ 的个数，
二维 Selberg 上界筛给出 $\Cnt_k(N)\le(8+o(1))\Sing(N)N/(\log N)^2$（参见 \cite{HR74}；常数 $8$ 以原文为准）；
又 $N\Haar(G_k^N)\sim4e^{-2\gamma}\Sing(N)N/(\log N)^2$（G1 注 rem:mertens），比值 $\le 8/(4e^{-2\gamma})=2e^{2\gamma}$。

(c) $N\Haar(G_k^N)=\Haar(U_k)^2(\log N)^2\cdot\frac{R_k(N)}{\Sing(N)}\cdot\Sing(N)\frac{N}{(\log N)^2}$；
第一个因子 $\to4e^{-2\gamma}$（Mertens 第三定理，$\prod_{p\le\sqrt N}(1-\frac1p)\sim2e^{-\gamma}/\log N$），
第二个因子 $=1+O(1/(\sqrt N\log N))$（G1 注 rem:mertensrate 的 $D$）。结合 (a) 即得等价。
\end{proof}

\begin{rem}[对任务表述“误差压倒主项”的修正]\label{rem:overwhelm}
在 $k=\pi(\sqrt N)$ 处被“压倒”的是 Legendre 型\emph{误差界}（命题 \ref{prop:top}(b)，超过主项 $e^{\Theta(\sqrt N/\log N)}$ 倍），
\emph{不是}真实误差：真实误差无条件地 $\le(2e^{2\gamma}-1+o(1))$ 倍主项（命题 \ref{prop:toperr}(b)），
且在 HL 猜想下恰为主项的 $1-e^{2\gamma}/4\approx20.7\%$（命题 \ref{prop:toperr}(c)）。
窗口太短的后果不是“误差无界”，而是“主项换了”：正确的主项不再是时间 $\to\infty$ 平均 $N\Haar(G_k^N)$，
而是它乘以一个依赖 $s$ 的因子（\S\ref{sec:mertens}）。另外，无条件地证明 $E_k/(N\Haar)>-1$（即 $\Cnt_k(N)>0$）
对充分大的 $N$ 就是二元哥德巴赫猜想本身，受奇偶性障碍阻挡，本稿不触及。
\end{rem}

%=====================================================================
\section{Mertens 因子 $4e^{-2\gamma}$ 的来源}\label{sec:mertens}

Buchstab 函数 $\omega:[1,\infty)\to\mathbb R$ 定义为 $\omega(u)=1/u$（$1\le u\le2$），$(u\,\omega(u))'=\omega(u-1)$（$u>2$）。
已知 $\omega(u)\to e^{-\gamma}$（$u\to\infty$），且 $\omega(u)-e^{-\gamma}$ 超指数衰减并无穷次变号（\cite[III.6]{Ten15}）。
例如 $\omega(2)=\frac12$，$\omega(3)=\frac{1+\log2}3=0.56438$，$e^{-\gamma}=0.56146$；第一次 $\omega(u)=e^{-\gamma}$ 约在 $u\approx2.48$。

\begin{prop}[一维 Buchstab 定理]\label{prop:buchstab}
令 $\Phi(x,y):=\#\{1\le n\le x:\ p\mid n\Rightarrow p>y\}$。对固定 $u>1$，$y=x^{1/u}$：
\[
\Phi(x,y)\sim\omega(u)\frac{x}{\log y},\qquad\text{从而}\qquad
\frac{\Phi(x,y)}{x\prod_{p\le y}(1-1/p)}\longrightarrow e^{\gamma}\omega(u).
\]
$u=2$ 时 $\Phi(x,\sqrt x)=\pi(x)-\pi(\sqrt x)+1$，比值 $\to e^\gamma/2=(2e^{-\gamma})^{-1}$：
Legendre/Haar 主项 $x\prod_{p\le\sqrt x}(1-\frac1p)$ 把 $\pi(x)$ 渐近地高估 $2e^{-\gamma}\approx1.1229$ 倍。
\end{prop}
\status{已证（经典：Buchstab \cite{Buc37}，de Bruijn \cite{dB50}；见 \cite[III.6]{Ten15}。第二式由 Mertens 第三定理）}

\begin{prop}[筛法基本引理给出的可控层]\label{prop:FL}
存在绝对常数 $s_0,C>0$，使得对偶数 $N$、$k\ge1$，记 $s:=\log(N/\log^4N)/\log p_k$，当 $s\ge s_0$ 时
\[
\rho_k(N)\le C\Bigl(e^{-s}+\frac1{\log N}\Bigr).
\]
因此对 $\varepsilon\gg1/\log N$，凡 $\log p_k\le c\,\log N/\log(1/\varepsilon)$ 的层都满足 $\rho_k(N)\le\varepsilon$；
这对应 $\log P_k\approx N^{c/\log(1/\varepsilon)}$，远超 Legendre 的 $\log P_k\approx\log N\log\log N/\log3$，但仍未到 $s=2$。
\end{prop}
\status{部分（引用筛法基本引理 \cite{HR74,FI10}，维数 $\kappa=2$；常数不显式）}
\begin{proof}[证明梗概]
取 $\mathcal A=\{0\le n<N\}$，$X=N$，$g(d)=\nu(d)/d$，$|\mathcal A_d|=Xg(d)+r_d$，$|r_d|\le\nu(d)\le\tau(d)$
（同命题 \ref{prop:legendre} 的论证）；$g(p)\le2/p$ 满足维数 $2$ 的条件。基本引理（例如 \cite[Ch.~2]{HR74}、\cite[\S6]{FI10}，
形式以原文为准）给出 $\Cnt_k(N)=N\Haar(G_k^N)(1+O(e^{-s}))+O\bigl(\sum_{d<D,\,d\mid P_k}|r_d|\bigr)$，$D=p_k^s$。
取 $D=N/\log^4N$，$\sum_{d<D}\tau(d)\le D(1+\log D)\le N/\log^3N$，除以 $N\Haar(G_k^N)\gg N/\log^2p_k$ 得 $O(1/\log N)$。
\end{proof}

\begin{conj}[二维 Buchstab]\label{conj:2dim}
设偶数 $N\to\infty$，$k=k(N)$ 满足 $\log N/\log p_k\to s\in[2,\infty)$。则
\[
\frac{\Cnt_k(N)}{N\Haar(G_k^N)}\longrightarrow\bigl(e^{\gamma}\omega(s)\bigr)^2,\qquad\text{即}\qquad
\frac{E_k(N)}{N\Haar(G_k^N)}\to\bigl(e^{\gamma}\omega(s)\bigr)^2-1 .
\]
\end{conj}
\status{猜想}
已知的部分：$s=2$ 时它等价于 HL 猜想（命题 \ref{prop:toperr}(c)，因 $(e^\gamma\omega(2))^2=e^{2\gamma}/4$）；
$s\ge s_0$ 时命题 \ref{prop:FL} 给出 $\limsup|\cdot-1|\ll e^{-s}$，与 $(e^\gamma\omega(s))^2=1+O(e^{-s})$ 相容但不蕴含极限。
\emph{启发式：}在 $p\le p_k$ 处的局部相关已全部计入 $\Haar(G_k^N)$；对 $p>p_k$ 假设 $n$ 与 $N-n$ 的“粗糙性”独立
（$\prod_{p>p_k}c_p(N)\to1$），并用一维 Buchstab 密度 $\omega(\log t/\log y)/\log y$ 描述 $t$ 附近的 $y$-粗糙数，得有限 $N$ 版本
\[
F_N(y):=e^{2\gamma}\frac1N\int_0^N\omega\Bigl(\frac{\log t}{\log y}\Bigr)\omega\Bigl(\frac{\log(N-t)}{\log y}\Bigr)dt
\qquad(\omega:=0\ \text{于}\ [0,1)),
\]
$N\to\infty$ 时 $F_N(N^{1/s})\to(e^\gamma\omega(s))^2$。极限曲线的若干值：
$(e^\gamma\omega(s))^2=0.7931,\,0.8630,\,0.8993,\,0.9556,\,1.0026,\,1.0104$（$s=2,\,2.1,\,2.165,\,2.3,\,2.5,\,3$）。
故（猜想地）$\rho_k>10\%$ 只发生在 $s\lesssim2.165$，即 $p_k\gtrsim N^{0.462}$；且 $E_k$ 在 $s\approx2.48$ 附近变号（$s$ 更大时 $E_k>0$）。

\begin{rem}[Mertens 因子的来源，及与 G1 注 rem:mertensrate 的衔接]\label{rem:mertenssource}
\begin{enumerate}[leftmargin=*,label=(\arabic*)]
\item \emph{来源。} 时间 $\to\infty$ 平均（Haar）给出主项 $N\Haar(G_k^N)$；它对 $k\le\kmax(N)$ 可严格控制（\S\ref{sec:legendre}），
在基本引理范围 $s\ge s_0$ 内仍正确到 $O(e^{-s})$（命题 \ref{prop:FL}）。但长度 $N$ 的窗口只在模 $\le N$ 的结构上等分布（注 \ref{rem:level}），
当 $s$ 降到 $2$ 时，窗口中被筛剩的数就是 $(\sqrt N,N)$ 中的素数，其密度是 $1/\log t$ 而不是 Haar 密度
$\prod_{p\le\sqrt N}(1-\frac1p)\approx2e^{-\gamma}/\log N$。每个坐标损失因子 $e^\gamma\omega(2)=e^\gamma/2$（一维是定理，命题 \ref{prop:buchstab}），
两个坐标（$n$ 与 $N-n$）合起来就是
\[
4e^{-2\gamma}=\bigl(e^\gamma\omega(2)\bigr)^{-2}.
\]
即：Mertens 因子不是“奇异级数错了”（$R_k(N)/\Sing(N)\to1$，G1 命题 prop:ratio），而是“窗口太短、Haar 平均不再是窗口平均”的定量表现。
其严格部分是命题 \ref{prop:toperr}(a)(c)；“每坐标一个 $e^\gamma\omega(2)$”在二维情形等价于 HL，属猜想 \ref{conj:2dim}。
\item \emph{收敛速度 $1/\log N$。} 在 $s=2$ 处 $\omega(\log t/\log\sqrt N)=\frac{\log N}{2\log t}$（$\sqrt N\le t\le N$），故
\[
F_N(\sqrt N)=\frac{e^{2\gamma}}4\cdot\frac{(\log N)^2}{N}\int_{\sqrt N}^{N-\sqrt N}\frac{dt}{\log t\,\log(N-t)}
=\frac{e^{2\gamma}}{4}\cdot\frac1{B}\,\bigl(1+O(N^{-1/2})\bigr),
\]
其中 $B=N/(L^2I(N))=1-\frac2L+O(L^{-2})$ 正是 G1 注 rem:mertensrate 中的因子 $B$（$L=\log N$）。
于是 G1 的分解 $\frac{N\Haar(G_k^N)}{r(N)}=A\cdot B\cdot D\cdot C$ 可读作：$A\to4e^{-2\gamma}$ 是“Buchstab 在 $s=2$ 的极限”，
$B$ 是同一 Buchstab 积分的有限 $N$ 修正（$1/\log N$ 量级，慢收敛的唯一主要来源），$D\to1$ 是奇异级数尾部，$C$ 是 HL 本身的误差。
\item \emph{数值（说明性，不作证明）。} G1 已报告 $N=10^8$ 时 $N\Haar(G_k^N)/r(N)=1.1197$（results/V9.json，字段
\texttt{N\_Haar\_G\_over\_r}），故 $1-r/(N\Haar)=0.1069$；D2 在同一 $N=10^8$、顶层 $k=1229$ 处给出
$\rho_k=0.1073$（results/D2.json，\texttt{cases[N=100000000].layers[k=1229].relative\_error}）。
两者之差 $0.0004$ 落在命题 \ref{prop:toperr}(a) 的边界项范围 $2\pi(10^4)/(N\Haar)\le2458/652589\approx0.0038$ 之内。
$0.107$ 距离极限 $0.207$ 只走了约一半，与 $B=1-2/L+\cdots$ 的慢收敛一致；渐近值 $0.207$ 是由 Mertens 定理推出的（在 HL 下），而不是由这些数推出的。
\end{enumerate}
\end{rem}

%=====================================================================
\section{与 Codex D2/D3 的对照}\label{sec:D2}

\paragraph{数据来源与字段含义。} results/D2.json 与 results/D3.json 在本稿写作期间由 Codex 生成（开始写作时尚不存在，
写作中途出现，文件时间戳 07:49/07:51），\texttt{status="passed"}。只读引用。
\begin{itemize}[leftmargin=*]
\item D2：\texttt{method} = ``direct sequential boolean sieve; count only newly removed positions; 0<n<N''（直接筛，不依赖 CRT 乘积公式）。
\texttt{cases[]} 的字段：\texttt{N}，\texttt{labels}（选样类别），\texttt{S}（$=\Sing(N)$），
\texttt{layers[]}：\texttt{k}、\texttt{p}（$=p_k$）、\texttt{count}（$=\Cnt_k(N)$）、\texttt{Haar}（$=\Haar(G_k^N)$）、
\texttt{E}（$=E_k(N)$）、\texttt{relative\_error}（$=\rho_k(N)$）、\texttt{x}（$=\log P_k/\log N$）；
\texttt{critical}：\texttt{relative\_error} 首次 $>0.1$ 的层（字段 \texttt{k,p,x,relative\_error}），无则为 \texttt{null}。
我在 g2\_scratch/ 中复核了 \texttt{relative\_error}$=|E|/(N\cdot\mathrm{Haar})$、\texttt{x}$=\log P_k/\log N$、\texttt{E}$=$\texttt{count}$-N\cdot$\texttt{Haar} 对全部层成立。
因此 D2 的 $k_c(N)$ 就是本稿 $\varepsilon=0.1$ 时“$\rho_k>\varepsilon$ 的第一层”。
\item D3：\texttt{ratio}$=r(N)/(\Sing(N)I(N))$，$I(N)=\int_2^{N-2}\frac{dt}{\log t\log(N-t)}$，即 G1 分解中的 $1/C$。
\end{itemize}

\paragraph{严格的相容性要求（由命题 \ref{prop:legendre} 推出，必须成立）。}
对每个 $N$：$k_c(N)>\kmax^{(N)}(N,0.1)\ge\kmax(N,0.1)$；并且每层 $|E_k(N)|\le\Tmix_k(N)$。
D2 的 24 个 $N$、9871 层全部满足（后者见 \S\ref{sec:legendre} 的注）。

\paragraph{对照表（说明性）。} 记 $s_c:=\log N/\log p_c$；“顶层”为 $k=\pi(\sqrt N)$；最后一列为各层 $\rho_k$ 的最大值及其所在的 $s$。
前三列来自本稿的定义计算，其余来自 results/D2.json 的上述字段。
\begin{center}
\small
\begin{tabular}{rcc|cccc|cc|c}
\hline
$N$ & $\kmax$ ($x_{\max}$) & $\kmax^{(N)}$ & $k_c$ & $p_c$ & $x_c$ & $s_c$ & 顶层 $k$ & 顶层 $\rho$ & $\max_k\rho_k$ ($s$)\\
\hline
$8192$ & 5 (0.86) & 5 & 17 & 59 & 5.4 & 2.21 & 24 & 0.0635 & 0.1111 (2.14)\\
$10000$ & 5 (0.84) & 6 & 16 & 53 & 4.9 & 2.32 & 25 & 0.0208 & 0.1079 (2.26)\\
$12468$ & 5 (0.82) & 5 & \multicolumn{4}{c|}{无（null）} & 29 & 0.0179 & 0.0429 (2.21)\\
$30030$ & 6 (1.00) & 8 & \multicolumn{4}{c|}{无（null）} & 40 & 0.0077 & 0.0523 (2.33)\\
$90090$ & 7 (1.15) & 9 & \multicolumn{4}{c|}{无（null）} & 62 & 0.0402 & 0.0873 (2.17)\\
$100000$ & 7 (1.14) & 8 & \multicolumn{4}{c|}{无（null）} & 65 & 0.0249 & 0.0708 (2.22)\\
$102468$ & 7 (1.14) & 7 & \multicolumn{4}{c|}{无（null）} & 66 & 0.0016 & 0.0584 (2.16)\\
$131072$ & 7 (1.12) & 7 & \multicolumn{4}{c|}{无（null）} & 72 & 0.0387 & 0.0932 (2.15)\\
$510510$ & 8 (1.22) & 11 & 82 & 421 & 30.0 & 2.18 & 127 & 0.0624 & 0.1025 (2.14)\\
$990990$ & 9 (1.39) & 11 & 98 & 521 & 35.7 & 2.21 & 167 & 0.0681 & 0.1077 (2.13)\\
$1000000$ & 9 (1.39) & 10 & 99 & 523 & 36.2 & 2.21 & 168 & 0.0674 & 0.1074 (2.15)\\
$1002468$ & 9 (1.39) & 9 & 100 & 541 & 36.6 & 2.20 & 168 & 0.0637 & 0.1077 (2.14)\\
$1021020$ & 9 (1.39) & 11 & 103 & 563 & 37.9 & 2.18 & 169 & 0.0681 & 0.1077 (2.12)\\
$1048576$ & 9 (1.39) & 9 & 111 & 607 & 41.5 & 2.16 & 172 & 0.0628 & 0.1047 (2.13)\\
$8388608$ & 11 (1.63) & 11 & 195 & 1187 & 71.8 & 2.25 & 418 & 0.0893 & 0.1260 (2.10)\\
$9999990$ & 11 (1.61) & 13 & 209 & 1289 & 77.2 & 2.25 & 446 & 0.0915 & 0.1245 (2.10)\\
$10000000$ & 11 (1.61) & 11 & 212 & 1301 & 78.6 & 2.25 & 446 & 0.0910 & 0.1248 (2.10)\\
$10002468$ & 11 (1.61) & 11 & 215 & 1319 & 79.9 & 2.24 & 446 & 0.0894 & 0.1235 (2.11)\\
$10210200$ & 11 (1.61) & 13 & 212 & 1301 & 78.5 & 2.25 & 452 & 0.0909 & 0.1255 (2.11)\\
$99999900$ & 13 (1.81) & 15 & 490 & 3511 & 185.5 & 2.26 & 1229 & 0.1079 & 0.1384 (2.09)\\
$100000000$ & 13 (1.81) & 14 & 487 & 3469 & 184.2 & 2.26 & 1229 & 0.1073 & 0.1382 (2.09)\\
$100002468$ & 13 (1.81) & 13 & 488 & 3491 & 184.6 & 2.26 & 1229 & 0.1078 & 0.1383 (2.09)\\
$100059960$ & 13 (1.81) & 15 & 492 & 3527 & 186.4 & 2.26 & 1229 & 0.1070 & 0.1378 (2.09)\\
$134217728$ & 13 (1.78) & 13 & 549 & 3967 & 208.5 & 2.26 & 1393 & 0.1097 & 0.1388 (2.09)\\
\hline
\end{tabular}
\end{center}

\paragraph{读法（说明性，均非证明）。}
\begin{enumerate}[leftmargin=*]
\item \emph{Legendre 界极为保守。} 在 $k=\kmax^{(N)}$ 处，D2 的实际 $\rho_k$ 为 $10^{-6}$--$10^{-3}$ 量级（例如 $N=10^8$、$k=14$：$1.0\times10^{-5}$），
$|E|/\Tmix_k$ 小于 $0.05$；混合界在 $N=10^8$ 附近从 $k\approx16$--$18$ 起就超过主项（命题 \ref{prop:kmax}(d) 的具体化），
而实际 $\rho_k$ 直到 $k\approx300$ 才超过 $5\%$。
\item \emph{$k_c$ 由 $s$ 决定，不由 Legendre 决定。} $k_c/\kmax$ 从约 $3$（$10^4$）增长到约 $37$（$10^8$）；
$x_c$ 从 $\approx5$ 增长到 $\approx185$--$209$，而 $x_{\max}\le1.81$。所有存在临界层的 $N$ 都有 $s_c\in[2.16,2.32]$，
$\rho_k$ 的最大值都出现在 $s\in[2.09,2.33]$。这与猜想 \ref{conj:2dim} 的图像一致：误差只在 $s$ 接近 $2$ 时变大。
按此图像 $k_c(N)\approx\pi(N^{1/s_c})$，阶为 $N^{0.44\text{--}0.46}/\log N$，与 Legendre 的 $\kmax\asymp\log N$ 完全不同量级；
命题 \ref{prop:FL} 严格地给出了中间的 $\log p_k\asymp\log N/\log(1/\varepsilon)$。
\item \emph{有限 $N$ 的 Buchstab 积分与 D2 的吻合。} 用 g2\_scratch/buchstab.py 数值计算启发式 $F_N(p_k)$（取 $y=p_k$），与 D2 的
$\Cnt_k/(N\Haar)=1+E/(N\cdot\mathrm{Haar})$ 比较：$N=10^8$ 时 $k=487,614,860,1044,1229$ 处分别为
$0.9000/0.9002$，$0.8784/0.8781$，$0.8620/0.8609$，$0.8689/0.8676$，$0.8927/0.8905$（D2/$F_N$）；
$N=10^6$ 时（$k=42,84,99,117,142,168$）差异在 $0.012$ 以内，$N=10^4$ 时差异达 $0.04$。有限 $N$ 下 $\rho_k$ 在 $s\approx2.09$ 达到峰值后向顶层回落，
是 $F_N$ 中 $1/\log N$ 修正（注 \ref{rem:mertenssource}(2)）在 $s$ 接近 $2$ 时最强所致；极限曲线本身在 $[2,2.48]$ 上单调。
D2 中 $E_k>0$ 的中间层（如 $N=10^6$，$s\approx3.2$；$N=10^8$，$s\approx2.7$--$3.4$）与 $\omega(s)>e^{-\gamma}$（$s\gtrsim2.48$）的符号预测一致。
这些吻合不证明猜想 \ref{conj:2dim}。
\item \emph{没有临界层的 $N$。} $N\in\{12468,30030,90090,100000,102468,131072\}$ 时各层 $\rho_k$ 均 $<10\%$。
若猜想 \ref{conj:2dim} 成立，则 $\max_k\rho_k\to1-(e^\gamma\omega(2))^2\approx0.207>0.1$，充分大的 $N$ 都会有 $k_c$；
但这不能从现有数据推出。$k_c$ 的存在与否在 $10^4$--$10^5$ 量级对 $N$ 的算术敏感（如 $\Sing(N)$ 较大的 $30030,90090$ 无临界层）。
\item \emph{D3：HL 误差 $C$ 很小。} results/D3.json 的 \texttt{high\_random.bins[0]}（$[10^7,10^8)$，2700 个分层均匀随机偶数）
给出 \texttt{mean}$=0.99962$、\texttt{std\_population}$=0.00154$；\texttt{full.bins} 中 $[10^6,10^7)$ 全部偶数 \texttt{mean}$=0.99887$、
\texttt{std\_population}$=0.00399$。即在此范围 $1/C=r/(\Sing I)$ 与 $1$ 相差约 $0.1\%$--$0.4\%$，
因此 D2 顶层约 $9\%$--$11\%$ 的偏差几乎全部来自 $A\cdot B$（Mertens 因子及其 $1/\log N$ 修正），而不是 HL 的失效。
这与注 \ref{rem:mertenssource} 的解释相容，但同样不构成证明。
\end{enumerate}

%=====================================================================
\section{小结与不确定之处}\label{sec:summary}

\begin{enumerate}[leftmargin=*]
\item \textbf{项数。} Legendre 展开按“（$d$，剩余类）”计有 $3^{k-\omega_k(N)}2^{\omega_k(N)}$ 项；“$2^{\omega(P_k)}$ 型误差”我们无法证明，
数值（$k\le6$）未见反例。本稿所有严格界都是 $\Tmix_k(N)\le2\cdot3^{k-2}$。
\item \textbf{$\kmax$。} $\kmax(N,\varepsilon)=\frac{\log(\varepsilon N)-2\log\log\log N}{\log3}+O(1)$，$x_{\max}\sim\log\log N/\log3$。
这是 \emph{Legendre} 可控的阶，不是真实可控的阶；后者由筛法参数 $s=\log N/\log p_k$ 决定（命题 \ref{prop:FL}、猜想 \ref{conj:2dim}）。
\item \textbf{显式常数。} 推论 \ref{cor:relative} 的 $c_*$ 依赖 \cite{RS62} 的不等式，本稿未重证，仅在 $y\le2\times10^5$ 抽查；
命题 \ref{prop:toperr}(b) 的常数 $8$ 与命题 \ref{prop:FL} 的 $s_0,C$ 均为引用，具体条件以原文为准。
\item \textbf{文献。} 本稿新增文献的元数据按记忆转录，\emph{未}在本稿中联网核对。
\item \textbf{启发式离散化。} $F_N(y)$ 用 $y=p_k$ 代入（$\Phi$ 的“$>y$”与 $G_k^N$ 的“$>p_k$”一致），$\omega$ 用步长 $10^{-4}$ 的梯形法数值求解；
它只是说明工具。
\item \textbf{“Polignac $=$ 时间 $\to\infty$ 平均”} 只对\emph{固定} $k$ 精确；真实的间隔 $N$ 素数计数同样需要 $p_k\approx\sqrt X$（注 \ref{rem:level}）。
\end{enumerate}

%=====================================================================
\section*{文献说明}
\cite{G1} 为本项目 G1 稿；\cite{HL23} 元数据沿用 G1（已核）。其余文献元数据未在本稿中联网核对，以原书/原文为准。

\begin{thebibliography}{99}
\bibitem{G1} （G1 工作稿）\emph{G1：哥德巴赫筛 $=$ 有限窗口的 Polignac 筛——字典、奇异级数与平均}，goldbach\_props.tex，2026-11-02。
\bibitem{HL23} G. H. Hardy, J. E. Littlewood, \emph{Some problems of `Partitio numerorum'; III: On the expression of a number as a sum of primes},
Acta Math. \textbf{44} (1923), 1--70. DOI: \url{10.1007/BF02403921}.
\bibitem{Buc37} A. A. Buchstab, \emph{Asymptotic estimates of a general number-theoretic function}（俄文），Mat. Sb. \textbf{2(44)} (1937), 1239--1246.
\bibitem{dB50} N. G. de Bruijn, \emph{On the number of uncancelled elements in the sieve of Eratosthenes},
Proc. Kon. Ned. Akad. Wetensch. \textbf{53} (1950), 803--812.
\bibitem{Ten15} G. Tenenbaum, \emph{Introduction to Analytic and Probabilistic Number Theory}, 3rd ed., Graduate Studies in Math.\ \textbf{163}, AMS, 2015（第 III.6 章：Buchstab 函数）。
\bibitem{HR74} H. Halberstam, H.-E. Richert, \emph{Sieve Methods}, Academic Press, London, 1974.
\bibitem{FI10} J. Friedlander, H. Iwaniec, \emph{Opera de Cribro}, AMS Colloquium Publications \textbf{57}, 2010.
\bibitem{RS62} J. B. Rosser, L. Schoenfeld, \emph{Approximate formulas for some functions of prime numbers}, Illinois J. Math. \textbf{6} (1962), 64--94.
\end{thebibliography}

\end{CJK}
\end{document}
